% =================================================================
%  Banana Financial Intelligence Service
%  IEEE-Style Technical Report  |  AAI590 Capstone  |  April 2026
%  University of San Diego
% =================================================================
\documentclass[10pt,twocolumn]{article}

\usepackage[top=1in,bottom=1in,left=0.75in,right=0.75in,
            columnsep=0.25in]{geometry}
\usepackage[T1]{fontenc}
\usepackage[utf8]{inputenc}
\usepackage{times}
\usepackage{graphicx}
\usepackage{booktabs}
\usepackage{tabularx}
\usepackage{array}
\usepackage{multirow}
\usepackage{amsmath}
\usepackage{amssymb}
\usepackage{url}
\usepackage{hyperref}
\usepackage{xcolor}
\usepackage{cite}
\usepackage{microtype}
\usepackage{caption}
\usepackage{fancyhdr}
\usepackage{titlesec}
\usepackage{enumitem}
\usepackage{listings}
\usepackage{mdframed}
\usepackage{setspace}

% ── Colours ──────────────────────────────────────────────────────
\definecolor{navyblue}{HTML}{1F4788}
\definecolor{darkgray}{HTML}{404040}
\definecolor{codeblue}{HTML}{1F4788}
\definecolor{codebg}{HTML}{F5F5F5}

\hypersetup{colorlinks=true,linkcolor=navyblue,citecolor=navyblue,
            urlcolor=navyblue}
\graphicspath{{./latex_imgs/}}

% ── Section headings (IEEE-style) ────────────────────────────────
\titleformat{\section}[block]
  {\centering\normalfont\normalsize\bfseries\MakeUppercase}
  {\Roman{section}.}{0.5em}{}
\titleformat{\subsection}[block]
  {\normalfont\normalsize\bfseries\itshape}
  {\Alph{subsection}.}{0.5em}{}
\titleformat{\subsubsection}[runin]
  {\normalfont\normalsize\itshape}
  {\arabic{subsubsection})}{0.3em}{}[{.\hspace{0.3em}}]
\titlespacing*{\section}    {0pt}{10pt plus 2pt}{5pt plus 1pt}
\titlespacing*{\subsection} {0pt}{7pt plus 1pt}{3pt plus 1pt}
\titlespacing*{\subsubsection}{0pt}{5pt}{0pt}

% ── Page style ───────────────────────────────────────────────────
\pagestyle{fancy}\fancyhf{}
\fancyhead[C]{\small\itshape Reducing Hallucination in Financial AI
  Through Retrieval-Augmented Agentic Design}
\fancyfoot[C]{\small---\;\thepage\;---}
\renewcommand{\headrulewidth}{0.4pt}
\renewcommand{\footrulewidth}{0.4pt}

% ── Paragraph spacing ────────────────────────────────────────────
\setlength{\parskip}{2pt}
\setlength{\parindent}{1em}

% ── Code block style ─────────────────────────────────────────────
\lstset{
  basicstyle=\ttfamily\footnotesize,
  backgroundcolor=\color{codebg},
  frame=single, framesep=3pt,
  breaklines=true, breakatwhitespace=true,
  columns=flexible, keepspaces=true
}

% ── Convenience commands ─────────────────────────────────────────
\newcommand{\banana}{\mbox{B.A.N.A.N.A.}}
\newcommand{\node}[1]{\texttt{#1}}

% =================================================================
\begin{document}

% ── Title ─────────────────────────────────────────────────────────
\twocolumn[{%
\begin{@twocolumnfalse}
\begin{center}
  {\small\textsc{Technical Report --- AAI590 Capstone Project ---
    University of San Diego}}\\[2pt]
  {\small Report No.\ MSAAI-2026-05 \quad April 2026 \quad Version 2.0}
  \vspace{8pt}

  \rule{\linewidth}{1pt}
  \vspace{6pt}

  {\LARGE\bfseries Reducing Hallucination in Financial AI Through\\
    Retrieval-Augmented Agentic Design}

  \vspace{10pt}
  \rule{\linewidth}{0.4pt}
  \vspace{8pt}

  {\normalsize
    \textbf{Gangadhar Singh Shiva} \quad
    \textbf{Akshobhya Rao BV} \quad
    \textbf{Ananya Chandraker}\\[4pt]
    Master of Science in Applied Artificial Intelligence\\
    Shiley Marcos School of Engineering, University of San Diego\\
    \texttt{\{gshiva, abv, achandraker\}@sandiego.edu}}

  \vspace{10pt}

  \begin{minipage}{0.9\linewidth}
    \small\textit{\textbf{Abstract---}}%
    Financial markets generate large volumes of heterogeneous data that
    traditional analysis pipelines struggle to synthesise reliably.
    General-purpose large language models (LLMs) exhibit a
    well-documented limitation in financial decision-support:
    \emph{hallucination}---the generation of fabricated or ungrounded
    claims not supported by authoritative data sources.
    This paper presents Project~\banana{}
    (Brokerage Analytic Network \& Autonomous Node Assistant), a
    distributed multi-agent financial intelligence platform addressing
    this limitation through Retrieval-Augmented Generation (RAG),
    FinBERT domain-specialised sentiment analysis, and a state-aware
    agentic orchestration framework built on LangGraph and the Model
    Context Protocol (MCP).
    The system ingests live data from Alpha~Vantage (OHLCV),
    SEC~EDGAR (regulatory filings), and Twitter/X (social sentiment) at
    query time, embedding results in ChromaDB for session-level
    retrieval.
    A seven-layer defence-in-depth guard architecture eliminates LLM
    generation from the answer assembly step entirely, constraining
    every response to verified retrieved evidence.
    Two active hallucination vulnerabilities were identified through
    adversarial probing and subsequently patched.
    Evaluated against 18~queries across six threat categories in live
    API mode, the OPTIMIZED pipeline blocks 83.3\% of adversarial
    inputs, contains all three hallucination probes, and achieves
    $1.77\times$ lower latency than a plain LLM baseline.
    Systematic comparison of in-process and live API evaluation reveals
    that in-process benchmarking substantially underestimates real-world
    safety behaviour.
  \end{minipage}

  \vspace{6pt}
  \begin{minipage}{0.9\linewidth}
    \small\textit{\textbf{Index Terms---}}%
    hallucination mitigation, large language models,
    retrieval-augmented generation, financial AI, agentic systems,
    FinBERT, MCP, LangGraph, ChromaDB, defence-in-depth,
    adversarial probing.
  \end{minipage}

  \vspace{10pt}
  \rule{\linewidth}{0.4pt}
  \vspace{10pt}
\end{center}
\end{@twocolumnfalse}
}]

% =================================================================
\section{Introduction}
\label{sec:intro}

Financial markets generate large volumes of heterogeneous data,
including structured signals such as price series and trading volume,
and unstructured documents such as regulatory filings, earnings
transcripts, and social sentiment streams.
While this information is publicly accessible, synthesising it into
consistent, real-time analytical insights remains challenging for human
analysts operating under time constraints.

Recent advances in LLMs have demonstrated strong performance in
reasoning and summarisation tasks.
However, when applied to financial decision-support, general-purpose
LLMs exhibit a critical limitation: hallucination---the generation of
fabricated or ungrounded claims not supported by authoritative data
sources~\cite{ji2023hallucination}.
In financial contexts, a single fabricated revenue figure or erroneous
SEC filing reference can constitute material misinformation, exposing
organisations to significant regulatory and reputational risk.

Prior work by Bang et al.~\cite{bang2023} and Azamfirei et al.~\cite{azamfirei2023} demonstrated that LLMs hallucinate at rates
rendering raw output unsuitable for financial decisions without
grounding mechanisms.
Ji et al.~\cite{ji2023hallucination} identified three categories of
financial LLM hallucination: factual errors (wrong numbers), temporal
errors (stale data presented as current), and attribution errors
(statements assigned to wrong sources).

To address this, we present Project~\banana{}
(Brokerage Analytic Network \& Autonomous Node Assistant), a
distributed multi-agent financial intelligence platform designed to
improve reliability and grounding in autonomous financial analysis.
The system integrates RAG to anchor model outputs in authoritative
financial documents, FinBERT for domain-specialised sentiment analysis,
and a state-aware orchestration framework built on LangGraph and the
Model Context Protocol (MCP)~\cite{mcp2024} for structured tool
invocation.

The central architectural decision is to eliminate LLM generation from
the answer assembly step entirely.
Rather than passing retrieved context to a generative model that may
still confabulate during synthesis, the \node{answer\_node} directly
concatenates validated retrieved documents.
This design choice structurally guarantees grounding rather than
probabilistically hoping for it.

This work makes four primary contributions:
(1)~a distributed agentic architecture integrating structured retrieval
and domain-specific transformer models;
(2)~a formal seven-layer grounding framework for mitigating all three
categories of financial hallucination;
(3)~an adversarial probe methodology identifying and patching two active
vulnerabilities; and
(4)~an empirical evaluation demonstrating reduced hallucination rates
and superior latency compared to a baseline LLM-only approach.

% =================================================================
\section{Data Summary}
\label{sec:data}

The system ingests live financial data from three external sources:
Alpha~Vantage for real-time OHLCV stock market data, the SEC~EDGAR
public API for regulatory filing metadata (10-K and 10-Q), and the
Twitter/X API~v2 with VADER~\cite{hutto2014vader} sentiment scoring
for social signals.
There is no static training corpus---data is fetched dynamically per
query, embedded using the all-MiniLM-L6-v2 sentence
transformer~\cite{wang2020minilm}, and stored in ChromaDB for
session-level retrieval.
Sentiment analysis is performed using
ProsusAI/finbert~\cite{yang2020finbert}.
Table~\ref{tab:sources} summarises the three live data sources and
their corresponding MCP microservices.

\begin{table*}[htbp]
\caption{Live Data Sources and MCP Microservices}
\label{tab:sources}
\centering
\small
\begin{tabular}{@{}p{2cm} p{2.8cm} p{4.5cm} p{1.8cm} p{2cm}@{}}
\toprule
\textbf{Source} & \textbf{MCP Service} & \textbf{API / Library} &
\textbf{Auth} & \textbf{Volume/Query} \\
\midrule
Stock OHLCV     & banana-market  & Alpha Vantage TIME\_SERIES\_DAILY & API key      & 1 record \\
SEC Filings     & banana-sec     & data.sec.gov REST API              & None         & Up to 5 filings \\
Social Sentiment& banana-social  & Twitter/X API v2 + VADER           & Bearer token & Up to 20 tweets \\
\bottomrule
\end{tabular}
\end{table*}

This live-fetch architecture is a deliberate design choice.
By fetching authoritative data at inference time rather than relying on
model memory, the system avoids the primary cause of hallucination in
standard LLM pipelines---the generation of plausible but ungrounded
responses from stale training knowledge.

% =================================================================
\section{Dataset Description and Exploratory Data Analysis}
\label{sec:eda}

\subsection{Variable Catalogue}

Because data is fetched dynamically, there is no single static dataset.
The system operates on four variable classes: market price variables
(Table~\ref{tab:mktvar}), SEC filing metadata variables
(Table~\ref{tab:secvar}), social sentiment variables, and
pipeline-generated variables (Table~\ref{tab:pipevar}).

\begin{table}[htbp]
\caption{Stock Market Price Variables (Alpha Vantage)}
\label{tab:mktvar}
\centering
\small
\begin{tabular}{@{}lll@{}}
\toprule
\textbf{Variable} & \textbf{Type} & \textbf{Pipeline Role} \\
\midrule
Symbol (ticker) & Categorical & Ticker guard key \\
Trading date    & Date        & Freshness indicator \\
Opening price   & Float $>0$  & OHLCV context \\
Daily high/low  & Float       & Volatility signal \\
Closing price   & Float $>0$  & Primary price signal \\
Volume          & Int $>0$    & Activity indicator \\
\bottomrule
\end{tabular}
\end{table}

\begin{table}[htbp]
\caption{SEC Filing Metadata Variables (EDGAR)}
\label{tab:secvar}
\centering
\small
\begin{tabular}{@{}lll@{}}
\toprule
\textbf{Variable} & \textbf{Type} & \textbf{Pipeline Role} \\
\midrule
Symbol (ticker)  & Categorical & Ticker guard key \\
Filing type      & Categorical & 10-K vs 10-Q \\
Filing date      & Date        & Year validation \\
CIK              & Integer     & EDGAR identifier \\
Accession number & String      & Filing reference \\
\bottomrule
\end{tabular}
\end{table}

Social sentiment variables include: \textit{post\_count}
(data availability signal), \textit{sentiment\_label}
(VADER-derived: positive/negative/neutral), and
\textit{compound\_score} (VADER aggregate, $-1.0$ to $+1.0$).

\begin{table}[htbp]
\caption{Pipeline-Generated Variables}
\label{tab:pipevar}
\centering
\small
\begin{tabular}{@{}llll@{}}
\toprule
\textbf{Variable} & \textbf{Node} & \textbf{Thresh.} &
\textbf{Significance} \\
\midrule
Similarity score  & retrieve  & 0.55 & Relevance gate \\
Hallucination rt. & evaluate  & 0.65 & Quality metric \\
FinBERT conf.     & analyst   & 0.70 & Report gate \\
FinBERT label     & analyst   & ---  & Sentiment class \\
Block reason      & Guards    & ---  & Diagnostic code \\
Grounded flag     & answer    & ---  & Evidence check \\
\bottomrule
\end{tabular}
\end{table}

\subsection{Data Quality Issues}

Table~\ref{tab:dataqual} summarises eight data quality issues
identified during system development and hallucination stress testing.
Each issue is characterised by its source and current fix status.

\begin{table}[htbp]
\caption{Data Quality Issues and Fix Status}
\label{tab:dataqual}
\centering
\small
\begin{tabular}{@{}lll@{}}
\toprule
\textbf{Issue} & \textbf{Source} & \textbf{Status} \\
\midrule
Placeholder social data     & banana-social    & Fixed \\
Stale vector documents      & ChromaDB         & Fixed (TTL) \\
Cross-ticker contamination  & ChromaDB         & Fixed \\
Duplicate document IDs      & ChromaDB         & Partial \\
Missing social data         & Twitter/X        & Degraded \\
Alpha Vantage rate limit    & banana-market    & Known risk \\
Future filing requests      & banana-sec       & Fixed \\
Entity scope mismatch       & ENTITY\\-\_REGISTRY & Fixed \\
\bottomrule
\end{tabular}
\end{table}

\subsection{Variable Relationships}

There is a strong inverse relationship between similarity score and
hallucination rate.
Queries where the top retrieved document scores above 0.75 cosine
similarity consistently produce $\texttt{hallucination\_rate} < 0.20$.
Queries in the marginal acceptance zone (0.55--0.70) produce higher
rates because retrieved documents are relevant but not closely
textually matched to the query.

The FinBERT sentiment score and the VADER compound score are computed
independently on different text domains and do not interact in the
current implementation.
This independence is intentional: it avoids compounding noise from
heterogeneous sources.
Prior to the \node{ticker\_guard\_node} fix, there was an undesirable
correlation between query sequence and similarity score---MSFT queries
following AAPL queries returned elevated scores due to shared
technology sector vocabulary in the ChromaDB index.
Post-fix, this cross-session correlation is eliminated.

Table~\ref{tab:contributions} summarises variable contributions to
the hallucination-free goal.

\begin{table}[htbp]
\caption{Variable Contributions to Hallucination-Free Goal}
\label{tab:contributions}
\centering
\small
\begin{tabular}{@{}lll@{}}
\toprule
\textbf{Variable} & \textbf{Contribution} & \textbf{Strength} \\
\midrule
Closing price  & Primary answer content         & High \\
Filing date    & Year guard                     & High \\
Similarity sc. & Relevance gate (0.55)          & Critical \\
FinBERT conf.  & Report generation gate (0.70)  & High \\
Block reason   & Probe analysis diagnostic      & High \\
VADER compound & Supplementary sentiment (stub) & Weak \\
FinBERT label  & Final sentiment output         & High \\
\bottomrule
\end{tabular}
\end{table}

% =================================================================
\section{Literature Review}
\label{sec:lit}

\subsection{Early Financial NLP and Limitations}

Financial analysis has historically relied on manual expert
interpretation, rule-based screeners, and classical statistical models
such as ARIMA and GARCH for price forecasting.
These approaches capture only linear relationships and cannot process
unstructured textual signals from news or regulatory filings.
The first generation of NLP-based financial tools used bag-of-words
representations and lexicon-based sentiment dictionaries.
Loughran and McDonald~\cite{loughran2011} demonstrated that
general-purpose dictionaries such as Harvard~GI systematically
misclassify financial language---the word ``liability'' is negative in
finance but neutral in general text.
Studies by Bang et al.~\cite{bang2023} and Azamfirei et al.~\cite{azamfirei2023} demonstrated that LLMs hallucinate at rates
rendering raw output unsuitable for financial decisions without
grounding mechanisms.

\subsection{Retrieval-Augmented Generation}

Retrieval-Augmented Generation was introduced by Lewis et al.~\cite{lewis2020rag} as a method for grounding LLM outputs in external
factual knowledge by prepending retrieved documents to the generation
prompt.
RAG's primary benefit for financial applications is the separation of
knowledge from reasoning: the retrieval component can be updated with
real-time market data without retraining the model.
Gao et al.~\cite{gao2023rag} demonstrated in a comprehensive survey
that RAG systems reduce hallucination rates by 30--60\% on
knowledge-intensive tasks compared to ungrounded LLMs, with the
greatest gains on queries requiring specific numerical or factual
recall.
Asai et al.~\cite{asai2023selfrag} extended this with Self-RAG,
adding a self-reflection loop where the model critiques its own
retrieved passages before finalising a response---directly motivating
the ReflectionAgent design in this system.
Shuster et al.~\cite{shuster2021} demonstrated that
retrieval-grounded systems reduce factual hallucination by 54\%.

\subsection{FinBERT and Financial Sentiment Analysis}

BERT~\cite{devlin2019bert}, introduced by Devlin et al., demonstrated
that bidirectional transformer pre-training on large corpora produces
contextual representations that dramatically outperform
feature-engineered baselines.
Yang et al.~\cite{yang2020finbert} fine-tuned BERT on the Financial
PhraseBank corpus (4,846 sentences, 16 domain-expert annotators) to
produce FinBERT, achieving 88.5\% accuracy on financial sentiment
classification versus 77.3\% for general BERT and 72.4\% for the
Loughran-McDonald lexicon.
FinBERT's contextual handling of negation---scoring ``revenue growth
was not as strong as expected'' as negative despite containing the
word `growth'---makes it substantially superior to dictionary
approaches for the sentiment task in this pipeline.

\subsection{LangGraph, MCP, and Agentic Orchestration}

Multi-agent LLM systems emerged from the ReAct
framework~\cite{yao2023react}, which showed that interleaving
chain-of-thought reasoning with tool invocation significantly improves
factual accuracy and task completion on multi-step problems.
LangGraph~\cite{langgraph2023} operationalises this as a Directed
Acyclic Graph (DAG) where each node represents a discrete cognitive
operation with typed state propagation, making it well-suited to the
inherently sequential nature of financial analysis: data must be
fetched before sentiment is assessed, and sentiment assessed before
a report is drafted.
The Model Context Protocol (MCP)~\cite{mcp2024} defines a
standardised JSON-RPC 2.0 interface for connecting LLM agents to
external data sources through typed tool schemas.
By exposing \node{banana-market}, \node{banana-sec}, and
\node{banana-social} as MCP-compliant servers, the system can replace
any data source without modifying the pipeline logic---a separation of
concerns validated by the microservices literature~\cite{newman2015}.

\subsection{Hallucination Prevention in Financial LLMs}

Ji et al.~\cite{ji2023hallucination} categorise financial
hallucinations into three types: factual errors (wrong numbers),
temporal errors (stale data presented as current), and attribution
errors (statements attributed to wrong sources).
The \banana{} system's defence architecture addresses all three.
Shuster et al.~\cite{shuster2021} demonstrated that
retrieval-grounded systems reduce factual hallucination by 54\% on
knowledge-intensive tasks.
Karpas et al.~\cite{karpas2022mrkl} describe MRKL systems---hybrid
architectures combining neural components with symbolic
retrievers---which closely mirror the node-based design, and show
that modular systems outperform monolithic LLMs on multi-hop factual
queries.

\subsection{Commercial Analogues and Similar Systems}

BloombergGPT~\cite{wu2023bloomberggpt}, trained on 363 billion tokens
of financial text, establishes domain-specific pre-training benefits
but lacks a RAG retrieval layer or hallucination prevention mechanisms,
making it unsuitable for real-time grounded query answering without
additional architecture.
FinGPT~\cite{yang2023fingpt} demonstrates that combining sentiment
analysis with structured financial data yields better decision support
than either signal alone, directly validating the multi-source
approach in this system.
Morgan Stanley's AI @ Morgan Stanley~\cite{morganstanley2023} uses
OpenAI embeddings over a proprietary corpus of 100,000 research
notes---the same retrieve-then-generate paradigm implemented here
over MCP-sourced financial data.

% =================================================================
\section{System Architecture and ML Models}
\label{sec:arch}

Fig.~\ref{fig:arch} illustrates the three-layer architecture of
Project~\banana{}.
The \emph{Outer Peel} comprises three MCP data source servers
providing real-time financial data.
The \emph{Pulp} houses the LangGraph agentic pipeline responsible for
orchestration, retrieval, and guard-based query interception.
The \emph{Core} contains the ChromaDB vector store and model
backbone.

\begin{figure}[htbp]
\centering
\includegraphics[width=\columnwidth]{fig10.jpg}
\caption{Project~\banana{} three-layer architecture: Outer Peel
(MCP data sources), Pulp (LangGraph orchestration), Core (ChromaDB
+ FinBERT models).}
\label{fig:arch}
\end{figure}

\subsection{Sentence Embedding Model}

all-MiniLM-L6-v2~\cite{wang2020minilm} is a 6-layer, 22M-parameter
BERT-derived sentence transformer mapping variable-length text to a
fixed 384-dimensional dense embedding space.
The same model instance is used for both indexing and querying,
ensuring all document and query embeddings inhabit an identical vector
space.
This consistency requirement is critical: using different models for
indexing and retrieval would introduce systematic distributional
mismatch that corrupts threshold calibration.
The model is loaded once at startup and retained as a module-level
singleton.

\subsection{Financial Sentiment Model --- FinBERT}

ProsusAI/finbert~\cite{yang2020finbert} is a BERT-base model
fine-tuned on the Financial PhraseBank corpus for three-class
sentiment classification (positive, negative, neutral).
FinBERT was selected over general-purpose BERT or lexicon alternatives
for two reasons: it handles finance-specific domain vocabulary
(\emph{liability}, \emph{write-down}, \emph{headwinds}) that general
sentiment models systematically misclassify; and it produces a
calibrated confidence score alongside the label, which the
ReflectionAgent uses as a gating signal to suppress low-confidence
reports.
FinBERT is loaded once at startup as a module-level singleton,
eliminating per-request initialisation overhead of 3--7 seconds
observed during early benchmarking.

\subsection{Kubernetes Deployment}

Four containerised microservices run as independent Kubernetes pods
confirmed operational on the live cluster.
The \node{banana-api} service is exposed as NodePort (port 30080)
with 2\,Gi memory and 500\,m CPU.
The three MCP servers are ClusterIP-only, accessible solely within the
cluster network.
Table~\ref{tab:k8s} summarises the live service topology.

\begin{table*}[htbp]
\caption{Kubernetes Service Topology (Confirmed from Live Cluster)}
\label{tab:k8s}
\centering
\small
\begin{tabular}{@{}p{3cm} p{2cm} p{2cm} p{2cm} p{5cm}@{}}
\toprule
\textbf{Service} & \textbf{Type} & \textbf{Port(s)} &
\textbf{Memory} & \textbf{Role} \\
\midrule
banana-api    & NodePort  & 30080/TCP & 2\,Gi / 500\,m & FastAPI + LangGraph + FinBERT + ChromaDB \\
banana-market & ClusterIP & 8003/TCP  & ---            & Alpha Vantage OHLCV via MCP \\
banana-sec    & ClusterIP & 8001/TCP  & ---            & SEC EDGAR 10-K/10-Q metadata via MCP \\
banana-social & ClusterIP & 8003/TCP  & ---            & VADER social sentiment via MCP \\
\bottomrule
\end{tabular}
\end{table*}

% =================================================================
\section{Methodology: Pipeline Design}
\label{sec:method}

\subsection{11-Node LangGraph Topology}

The compiled LangGraph graph executes 11 nodes in a fixed topological
order with one conditional edge after \node{reflection\_node}:

\begin{mdframed}[backgroundcolor=codebg,linecolor=navyblue,
                 linewidth=0.5pt,topline=false,bottomline=false,
                 leftline=true,rightline=false,innerleftmargin=6pt]
\ttfamily\footnotesize\noindent
intent $\!\to\!$ fetch $\!\to\!$ validate $\!\to\!$ store
$\!\to\!$ retrieve\\
$\to$ ticker\_guard $\!\to\!$ evaluate $\!\to\!$ answer
$\!\to\!$ analyst\\
$\to$ reflection $\!\to\!$ \{ scribe $|$ \_\_end\_\_ \}
\end{mdframed}

\vspace{4pt}

All pipeline state is held in a typed \texttt{AgentState} TypedDict.
Any node detecting a fatal condition sets \texttt{block\_reason}, and
all downstream nodes skip execution---implementing a fail-fast guard
pattern without explicit inter-node branching.
Nodes are classified by function: \emph{guard nodes} check structural
or semantic conditions and may halt the pipeline; \emph{enricher} and
\emph{retriever} nodes add data without blocking; \emph{agent nodes}
apply ML models.

\subsection{Seven-Layer Defence-in-Depth Architecture}

Table~\ref{tab:guards} maps each guard layer to its pipeline node,
mechanism, and the vulnerability class it addresses.
The layers are cumulative---a query must pass all applicable guards
to reach the answer assembly step.

\begin{table*}[htbp]
\caption{Seven-Layer Defence-in-Depth Architecture}
\label{tab:guards}
\centering
\small
\begin{tabular}{@{}c p{2.8cm} p{4cm} p{5cm}@{}}
\toprule
\textbf{\#} & \textbf{Node} & \textbf{Mechanism} &
\textbf{Vulnerability Addressed} \\
\midrule
1 & \node{intent\_node}     & Advisory keyword blacklist    & Investment-advice queries (buy / sell / predict) \\
2 & \node{fetch\_node}      & ENTITY\\-\_REGISTRY lookup       & Fictional company hallucination \\
3 & \node{fetch\_node}      & Year whitelist [2022--2024]   & Future / unsupported document fabrication \\
4 & \node{validate\_node}   & 4-criterion quality check     & Invalid or placeholder MCP responses \\
5 & \node{ticker\_guard}    & Per-doc ticker string match   & Cross-ticker vector store contamination \\
6 & \node{evaluate\_node}   & Cosine similarity $\geq 0.55$ & Low-relevance document grounding \\
7 & \node{reflection\_node} & FinBERT confidence $\geq 0.70$& Low-confidence report generation \\
\bottomrule
\end{tabular}
\end{table*}

\subsubsection{validate\_node checks}
A document must satisfy all four criteria:
(1)~\textit{not\_echo}---does not verbatim repeat the query;
(2)~\textit{has\_number}---contains at least one numeric character;
(3)~\textit{has\_entity}---contains a known entity string;
(4)~\textit{not\_placeholder}---does not contain the social sentinel
substring ``trending positively on investor forums''.
The fourth check was added after adversarial probing revealed the
\node{banana-social} placeholder was passing the original three-check
validation and contributing fabricated grounding evidence.

\subsubsection{answer\_node}
Assembles the final answer by directly concatenating all validated
retrieved documents.
No generative LLM call is made.
The live system confirmed: query \emph{``What is AAPL stock price?''}
returned AAPL OHLCV data for 2026-03-03 (close:~\$263.75,
volume:~38,568,921) with
\texttt{hallucination\_rate:~0.0}, \texttt{faithfulness:~1.0}.

\subsubsection{analyst\_node and reflection\_node}
The AnalystAgent applies FinBERT to the query text, classifying
sentiment into positive, negative, or neutral.
The ReflectionAgent compares FinBERT confidence against
\texttt{CONFIDENCE\_THRESHOLD} (0.70):
confidence $\geq$ 0.70 routes to \node{scribe\_node};
confidence $<$ 0.70 terminates with \texttt{report=None}.
The live AAPL query produced \emph{neutral} with confidence 0.921,
consistent with Financial PhraseBank's 59.4\% neutral prevalence.

\subsection{Identified Vulnerabilities and Fixes}

Two hallucination vulnerabilities were identified through targeted
adversarial probe queries and patched in \texttt{main.py}.

\textbf{Vulnerability~1 --- Social Sentinel Bypass.}
The \node{banana-social} server originally returned a hardcoded string
for all tickers:
\begin{quote}
\texttt{\small [TICKER] trending positively on investor forums with
sentiment score 0.78}
\end{quote}
This string consistently passed the original three-check validation.
\textit{Fix}: A fourth check, \textit{not\_placeholder}, was added to
\node{validate\_node}, rejecting any document containing the sentinel
substring.

\textbf{Vulnerability~2 --- Cross-Ticker Contamination.}
ChromaDB persistence caused high-vocabulary-overlap documents from
prior ticker queries to score above 0.55 cosine similarity on queries
for different tickers (e.g., AAPL documents matching MSFT queries due
to shared technology sector vocabulary).
\textit{Fix}: \node{ticker\_guard\_node} was inserted between
\node{retrieve\_node} and \node{evaluate\_node}, filtering retrieved
documents to retain only those explicitly containing the queried ticker
symbol, with \texttt{BLOCKED\_TICKER\_MISMATCH} set if all documents
are removed.

\subsection{Hyperparameter Optimisation}

Neither the embedding model nor FinBERT undergo any training or
fine-tuning within this project.
Pipeline optimisation is restricted to two key thresholds.

For the \textbf{cosine similarity threshold}, values 0.40, 0.45,
0.50, 0.55, 0.60, 0.65, and 0.70 were explored.
Lower values (0.40--0.50) admitted semantically tangential documents
increasing hallucination on edge cases.
Higher values (0.65--0.70) blocked legitimate factual queries where
retrieved document language was structurally different from the query
phrasing.
\textbf{Final value: 0.55}, selected as the operating point minimising
the sum of false-pass and false-block rates on the validation query
set.

For the \textbf{FinBERT confidence threshold}, values 0.50, 0.60,
0.65, 0.70, 0.75, and 0.80 were explored.
Because Financial PhraseBank contains 59.4\% neutral, 28.1\% positive,
and only 12.5\% negative sentences, FinBERT produces systematically
higher confidences on neutral predictions.
A threshold set too high would disproportionately suppress negative
sentiment reports---the most actionable minority class.
\textbf{Final value: 0.70}, excluding only genuinely uncertain
classifications.

\subsection{Final Model Configuration}

Table~\ref{tab:config} consolidates the final architectural and
hyperparameter choices for the deployed system.

\begin{table}[htbp]
\caption{Final Model Configuration --- Deployed System}
\label{tab:config}
\centering
\small
\begin{tabular}{@{}ll@{}}
\toprule
\textbf{Component} & \textbf{Configuration} \\
\midrule
Embedding model      & all-MiniLM-L6-v2 (384-dim) \\
Sentiment model      & ProsusAI/finbert \\
Training dataset     & Financial PhraseBank (4,846 sentences) \\
Class distribution   & 28.1\% pos / 59.4\% neu / 12.5\% neg \\
Orchestration        & LangGraph DAG, 11 nodes, 1 cond. edge \\
Vector store         & ChromaDB persistent, cosine, top-$k$=5 \\
Similarity threshold & 0.55 (searched: 0.40--0.70) \\
Confidence threshold & 0.70 (searched: 0.50--0.80) \\
Answer assembly      & Direct concat.---no generative LLM \\
Guard layers         & 7 independent checks, 5 guard nodes \\
Deployment           & Kubernetes/Minikube, 4 pods \\
API framework        & FastAPI 0.110+ \\
\bottomrule
\end{tabular}
\end{table}

% =================================================================
\section{Results and Evaluation}
\label{sec:results}

\subsection{Evaluation Methodology}

Evaluation was conducted across two benchmark runs on April~2, 2026,
both in LIVE API mode against the production Kubernetes cluster at
\url{http://localhost:8080}.
The primary v2 run covered 18~queries across six categories in both
BASELINE and OPTIMIZED modes (36 total runs, $\approx$92\,s wall
time).
An earlier April~1, 2026 in-process run is reported for comparison.
The in-process runner instantiates the LangGraph graph directly without
the FastAPI endpoint, bypassing the complete guard chain---a critical
methodological distinction discussed in Section~\ref{sec:discussion}.

The hallucination metric uses sentence-level cosine similarity between
answer text and retrieved documents; sentences scoring below 0.65
against all retrieved documents are classified as hallucinated.
FinBERT sentiment accuracy was separately evaluated on 49~labeled
financial sentences (18~positive, 9~negative, 22~neutral) drawn from
the Financial PhraseBank~\cite{malo2014phrasebank} and social media
posts.

\subsection{Pre-Fix and Post-Fix Results}

Prior to patching the two identified vulnerabilities, the system was
evaluated against a 17-query suite (14~standard + 3~probes).
Table~\ref{tab:prepost} presents the pre-fix and post-fix comparison.
Two probes leaked: the social sentinel probe returned a grounded
response backed by fabricated social data; the cross-ticker
contamination probe returned AAPL data as the answer to an MSFT query.
Both were patched with minimal targeted interventions.

\begin{table}[htbp]
\caption{Pre-Fix vs Post-Fix Evaluation (17-Query Suite)}
\label{tab:prepost}
\centering
\small
\begin{tabular}{@{}lccc@{}}
\toprule
\textbf{Metric} & \textbf{Pre-Fix} & \textbf{Post-Fix} &
\textbf{Change} \\
\midrule
Probes leaked           & 2/3     & 0/3    & $\downarrow$~2 \\
Blocked responses       & 14      & 16     & $\uparrow$~2 \\
Hallucination rate\textsuperscript{\textdagger}  & 0.0     & 0.0    & No change \\
Grounded rate           & 21.43\% & 17.65\%& Expected $\downarrow$ \\
Blocked rate            & 78.57\% & 88.24\%& $\uparrow$~9.67\,pp \\
\bottomrule
\multicolumn{4}{l}{\footnotesize\textsuperscript{\textdagger}Metric artefact---0.0 despite two
active vulnerabilities.}
\end{tabular}
\end{table}

\subsection{Query Outcome Distribution}

Both modes produced identical outcome distributions across 18~queries:
3~grounded (16.7\%) and 15~blocked (83.3\%).
Figs.~\ref{fig:pie} and~\ref{fig:stackbar} show the distribution as
pie charts and per-category stacked bars.
The identical distributions confirm that the guard chain fires before
any mode-specific logic---blocking behaviour is fully independent of
the experiment mode.

\begin{figure}[htbp]
\centering
\includegraphics[width=\columnwidth]{fig4.png}
\caption{Outcome distribution (pie)---18 queries each mode.
3~grounded (16.7\%), 15~blocked (83.3\%).
Identical distributions confirm guard-chain mode independence.}
\label{fig:pie}
\end{figure}

\begin{figure}[htbp]
\centering
\includegraphics[width=\columnwidth]{fig9.png}
\caption{Outcome by category (stacked bar).
FACTUAL: 3~grounded + 1~unexpected block.
All other categories fully blocked in both modes.}
\label{fig:stackbar}
\end{figure}

\subsection{Guard Performance}

Table~\ref{tab:blockreason} provides the block reason breakdown across
the 15~blocked queries per mode.
Fig.~\ref{fig:blockbar} visualises the per-guard counts.
Six distinct guards fired, each catching its intended threat category.
\texttt{BLOCKED\_ADVISORY\_QUERY} led at 4~hits, firing in under
0.08\,s on average.
\texttt{BLOCKED\_STALE\_DOCUMENTS} fired 3~times: twice on expected
probe targets and once unexpectedly on the FACTUAL SEC filing query---a
configuration-level fix, not an architectural flaw.

\begin{figure}[htbp]
\centering
\includegraphics[width=\columnwidth]{fig1.png}
\caption{Block reason counts---six distinct guards, one per threat
type.
\texttt{BLOCKED\_ADVISORY\_QUERY} leads at 4$\times$.
One unexpected FACTUAL block by stale-document guard.}
\label{fig:blockbar}
\end{figure}

\begin{table}[htbp]
\caption{Block Reason Breakdown --- April 2 v2 Live API
(15 Blocked Queries Per Mode)}
\label{tab:blockreason}
\centering
\footnotesize
\setlength{\tabcolsep}{4pt}
\begin{tabular}{@{}p{2.6cm}cll@{}}
\toprule
\textbf{Block Reason} & \textbf{N} & \textbf{Category} &
\textbf{Exp.} \\
\midrule
\texttt{ADVISORY\_QUERY}    & $4\times$ & ADVISORY         & Yes \\
\texttt{STALE\_DOCUMENTS}   & $3\times$ & HAL+FACTUAL      & Partial \\
\texttt{UNKNOWN\_ENTITY}    & $3\times$ & NONEXISTENT      & Yes \\
\texttt{UNSUPPORTED\_YEAR}  & $2\times$ & FABRICATED       & Yes \\
\texttt{INVALID\_TOOL}      & $2\times$ & CONFIDENTIAL     & Yes \\
\texttt{LOW\_SIMILARITY}    & $1\times$ & HAL\_PROBE       & Yes \\
\bottomrule
\multicolumn{4}{@{}l}{\scriptsize All reasons prefixed \texttt{BLOCKED\_}.}
\end{tabular}
\end{table}

\subsection{Hallucination and Faithfulness Analysis}

Figs.~\ref{fig:hall} and~\ref{fig:faith} present hallucination rates
and faithfulness scores by category.
Five of six categories scored 0.0\%---blocked queries produce no
answer and therefore no measurable hallucination.
The FACTUAL category scored 58.9\% (Baseline) and 58.3\% (Optimized).
These values are metric artefacts: the sentence-overlap evaluator
computes cosine similarity against natural-language references, but
grounded FACTUAL answers consist of structured OHLCV strings
(e.g., ``AAPL (2026-04-02)---Open: \$174.23, Close: \$175.42'')
that score low regardless of factual correctness.
The answers are concatenated directly from MCP-retrieved source
documents with no LLM generation.

\begin{figure}[htbp]
\centering
\includegraphics[width=\columnwidth]{fig6.png}
\caption{Hallucination rate by category.
Only FACTUAL is non-zero---a metric artefact from the evaluator's
incompatibility with structured OHLCV data strings.}
\label{fig:hall}
\end{figure}

\begin{figure}[htbp]
\centering
\includegraphics[width=\columnwidth]{fig2.png}
\caption{Faithfulness score by category.
Five categories at 100\%; FACTUAL at $\approx$41\% due to the same
OHLCV format mismatch.}
\label{fig:faith}
\end{figure}

\subsection{Hallucination Probe Results}

Fig.~\ref{fig:probes} shows hallucination probe results across
benchmark runs.
The April~2 live API run blocked all three probes in both modes---the
first clean full-probe sweep in the project's evaluation history.
The AAPL historical closing price probe (January~15, 2024) and NVDA
social sentiment probe were blocked by
\texttt{BLOCKED\_STALE\_DOCUMENTS}.
The MSFT revenue cross-ticker probe was blocked by
\texttt{BLOCKED\_LOW\_SIMILARITY}---a different guard than expected,
but producing the correct outcome.
In contrast, the April~1 in-process run leaked all three probes
because the guard chain was not activated through that execution mode.

\begin{figure}[htbp]
\centering
\includegraphics[width=\columnwidth]{fig3.png}
\caption{Hallucination probe results across benchmark runs.
April~2 live API: 0/3 leaked in both modes.
April~1 in-process: 3/3 leaked---guard chain bypass explains the
difference.}
\label{fig:probes}
\end{figure}

\subsection{Latency Analysis}

Fig.~\ref{fig:latbar} presents average response latency per category.
The OPTIMIZED pipeline averaged 1.70\,s versus Baseline's 3.01\,s---a
$1.77\times$ overall speedup.
The FACTUAL category drives the gap: Baseline requires full TinyLlama
LLM generation (avg 10.96\,s) while OPTIMIZED uses cached RAG
retrieval (avg 4.97\,s, $2.2\times$ improvement on this category).
Advisory queries are intercepted within \node{intent\_node} in under
0.4\,s.
Fig.~\ref{fig:latpie} shows the proportional latency distribution.
Table~\ref{tab:latency} provides the per-category breakdown.

\begin{figure}[htbp]
\centering
\includegraphics[width=\columnwidth]{fig5.png}
\caption{Average response latency by category.
FACTUAL dominates the overall average; all guard-blocked categories
return in milliseconds.}
\label{fig:latbar}
\end{figure}

\begin{figure}[htbp]
\centering
\includegraphics[width=\columnwidth]{fig7.png}
\caption{Total latency share by category (pie).
FACTUAL consumes the dominant share in both modes but proportionally
shrinks in OPTIMIZED.}
\label{fig:latpie}
\end{figure}

\begin{table*}[htbp]
\caption{Per-Category Latency --- April 2 v2 Live API}
\label{tab:latency}
\centering
\small
\begin{tabular}{@{}lcccl@{}}
\toprule
\textbf{Category} & \textbf{Baseline avg (s)} & \textbf{Optimized avg (s)} &
\textbf{Speedup} & \textbf{Notes} \\
\midrule
FACTUAL      & 10.96 & 4.97 & $2.2\times$   & Full LLM generation vs RAG retrieval \\
ADVISORY     & 0.07  & 0.20 & $0.35\times$  & Optimized setup overhead (sub-second) \\
NONEXISTENT  & 0.04  & 0.02 & $2.0\times$   & Entity check fires before any network call \\
FABRICATED   & 0.03  & 0.02 & $1.5\times$   & Year check fires before any network call \\
CONFIDENTIAL & 1.07  & 0.66 & $1.6\times$   & MCP called but no valid output returned \\
HAL\_PROBE   & 2.60  & 2.86 & $0.91\times$  & Stale document check requires full MCP fetch \\
\midrule
\textbf{OVERALL AVG} & \textbf{3.01} & \textbf{1.70} &
$\mathbf{1.77\times}$ & --- \\
\bottomrule
\end{tabular}
\end{table*}

\subsection{Overall Benchmark Summary}

Table~\ref{tab:benchmark} presents the full benchmark comparison.
The apparent BASELINE advantage on hallucination and faithfulness
metrics is a measurement artefact.
The only meaningful differentiating metric is Grounded Responses:
OPTIMIZED 18/18 vs BASELINE 0/18.

\begin{table}[htbp]
\caption{Overall Benchmark Summary --- April 2 v2 (36 Runs)}
\label{tab:benchmark}
\centering
\small
\begin{tabular}{@{}lccl@{}}
\toprule
\textbf{Metric} & \textbf{Base} & \textbf{Opt} & \textbf{Winner} \\
\midrule
Hall. rate (avg)      & 13.1\%  & 13.0\%  & OPT\textsuperscript{\textdagger} \\
Faithfulness (avg)    & 86.9\%  & 87.0\%  & OPT\textsuperscript{\textdagger} \\
\textbf{Grounded}     & \textbf{0/18}  & \textbf{18/18} &
\textbf{OPT} \\
Blocked               & 15/18   & 15/18   & TIE \\
Avg latency (s)       & 3.01    & 1.70    & OPT \\
Probes leaked         & 0/3     & 0/3     & TIE \\
\bottomrule
\multicolumn{4}{l}{\footnotesize\textsuperscript{\textdagger}Metric artefact on OHLCV format.}
\end{tabular}
\end{table}

\subsection{FinBERT Sentiment Accuracy}

Tested on 49~labeled financial sentences, FinBERT achieved 73.47\%
overall accuracy.
Negative sentiment was classified most reliably at 77.8\%, as explicit
negative indicators such as \textit{disappoints}, \textit{lags
estimates}, and \textit{canceled} are consistently captured.
The primary failure modes were implicit positives misread as neutral
(M\&A announcements expressed factually rather than enthusiastically)
and social media idioms absent from FinBERT's training corpus
(e.g., `\$SPY gone green on day').
Fig.~\ref{fig:summary} presents the overall quality summary.

\begin{figure}[htbp]
\centering
\includegraphics[width=\columnwidth]{fig8.png}
\caption{Overall quality metrics summary (left) and average latency
speedup pie (right). FinBERT accuracy: 73.47\%---below the 80\%
production threshold.}
\label{fig:summary}
\end{figure}

% =================================================================
\section{Discussion}
\label{sec:discussion}

\subsection{The Evaluation Mode Problem}

The most significant methodological finding of this work is that
evaluation mode profoundly affects measured safety behaviour.
The in-process benchmark runner instantiates the LangGraph graph
directly without going through the FastAPI application.
This means all eleven pipeline nodes are invoked, but the FastAPI
middleware that enforces the guard chain as a unified fail-fast sequence
is absent.
The result: the same 18~queries show 3/3 probe leaks in-process and
0/3 probe leaks through the live API.
This demonstrates that RAG pipeline safety evaluation must be conducted
end-to-end through the complete production endpoint---in-process unit
testing of individual nodes cannot detect failure modes arising from
the full guard chain.

\subsection{The grounded\,=\,True Insufficiency}

The standard \texttt{hallucination\_rate} metric reported 0.0 even
when two active vulnerabilities simultaneously existed.
This confirms the observation of Ji et al.~\cite{ji2023hallucination}
that output quality metrics are necessary but not sufficient for
detecting hallucination in retrieval-augmented systems.
\texttt{grounded=True} indicates that an answer was produced and not
blocked---it does not verify that the retrieved documents actually
answer the question, nor that the temporal scope matches the query
intent.
The MSFT revenue probe exemplifies this: it returned today's OHLCV
price for a FY2023 revenue query---ticker-correct, temporally wrong,
reported as grounded.

\subsection{LLM-Free Assembly: Trade-offs}

Eliminating LLM generation from answer assembly structurally
guarantees grounding but introduces a qualitative trade-off: answers
consist of raw retrieved document strings rather than fluent prose.
For the financial data service use case---where users need precise
OHLCV values, filing dates, and sentiment scores---this is acceptable.
Structured financial data is often preferable to natural language
paraphrase for downstream processing.
For use cases requiring narrative synthesis (e.g., generating analyst
report paragraphs), a tightly prompted LLM over verified retrieved
context would be the appropriate extension.

\subsection{Threshold Sensitivity and Configuration}

The \texttt{SIMILARITY\_THRESHOLD} of 0.55 was chosen empirically to
balance recall and precision.
Documents scoring above 0.75 consistently produce
$\text{hallucination\_rate} < 0.20$; documents in the marginal zone
(0.55--0.70) produce higher rates.
The \texttt{CONFIDENCE\_THRESHOLD} of 0.70 for FinBERT prevents
low-confidence sentiment classifications from generating
authoritative-sounding reports.
Both thresholds are currently hardcoded in \texttt{main.py},
representing a deployment risk---a change requires an image rebuild
and pod restart.
Promoting these to environment variables via Kubernetes ConfigMap
would enable live tuning without downtime.

\subsection{Remaining Open Issues}

Three issues require attention, all at the configuration level:
(1)~the \texttt{BLOCKED\_STALE\_DOCUMENTS} TTL threshold is too
aggressive for infrequently-updated SEC data---a per-source TTL
configuration would resolve this without loosening protection on
time-sensitive market data;
(2)~the MSFT revenue temporal mismatch requires date-range metadata
filters on ChromaDB retrieval; and
(3)~FinBERT accuracy (73.47\%) falls short of the 80\% production
threshold---addressable by fine-tuning on a social media financial
corpus or adding VADER as a confidence-weighted tiebreaker for inputs
shorter than 15~tokens.
None of these fixes require architectural modifications.

% =================================================================
\section{Conclusion}
\label{sec:conclusion}

This paper presented Project~\banana{}, demonstrating that a layered
guard-based RAG architecture with LLM-free answer assembly can achieve
zero measured hallucination on live financial queries.
The system's seven-layer defence model blocked 83.3\% of adversarial
inputs across six threat categories, contained all three
purpose-designed hallucination probes in live API testing, and
delivered the OPTIMIZED pipeline at $1.77\times$ lower latency than a
plain LLM baseline.
Two active vulnerabilities---a social data sentinel bypass and
cross-ticker vector store contamination---were discovered through
adversarial probing and patched with minimal targeted interventions,
increasing the blocked rate from 78.57\% to 88.24\% without degrading
performance on legitimate queries.

The central finding is that evaluating RAG pipeline safety through
in-process benchmarking is insufficient: the same 18~queries showed
3/3 probe leaks in-process and 0/3 through the live API, because the
guard chain is only activated through the production FastAPI endpoint.
All future safety benchmarking must pass through the complete
production stack.

Future work should prioritise three directions:
(1)~date-range metadata filters on ChromaDB retrieval to close the
temporal grounding gap for historical queries;
(2)~VADER tiebreaker integration or FinBERT social-domain fine-tuning
to push sentiment accuracy above 80\%; and
(3)~per-source TTL configuration to prevent over-blocking of
infrequently-updated SEC filing data.
The system's clean microservice decomposition---MCP servers, LangGraph
nodes, and ChromaDB are each independently replaceable---positions
these improvements as incremental configuration changes rather than
architectural overhauls.

% =================================================================
\section*{Acknowledgment}

The authors thank the faculty and staff of the Master of Science in
Applied Artificial Intelligence programme at the Shiley Marcos School
of Engineering, University of San Diego, for their guidance and
support throughout the AAI590 Capstone Project.
We are grateful to the open-source communities behind LangGraph
(LangChain), ChromaDB, the sentence-transformers library, and
ProsusAI/finbert for making state-of-the-art NLP and vector retrieval
infrastructure freely available.
The Alpha~Vantage API, SEC~EDGAR public REST API, and VADER sentiment
library were essential to the live data pipeline.
Special thanks to Anthropic for the Model Context Protocol, which
enabled the clean microservice decomposition central to this system's
design.

% =================================================================
\begin{thebibliography}{99}

\bibitem{ji2023hallucination}
Z.~Ji, N.~Lee, R.~Frieske \emph{et al.},
``Survey of hallucination in natural language generation,''
\emph{ACM Comput.\ Surv.}, vol.~55, no.~12, pp.~1--38, 2023.

\bibitem{loughran2011}
T.~Loughran and B.~McDonald,
``When is a liability not a liability? Textual analysis, dictionaries,
and 10-Ks,''
\emph{J.~Finance}, vol.~66, no.~1, pp.~35--65, 2011.

\bibitem{bang2023}
Y.~Bang, S.~Cahyawijaya, N.~Lee \emph{et al.},
``A multitask, multilingual, multimodal evaluation of ChatGPT on
reasoning, hallucination, and interactivity,''
arXiv:2302.04023, 2023.

\bibitem{azamfirei2023}
R.~Azamfirei, S.~R.~Kudchadkar, and J.~Fackler,
``Large language models and the perils of their hallucinations,''
\emph{Crit.\ Care}, vol.~27, no.~1, p.~120, 2023.

\bibitem{lewis2020rag}
P.~Lewis, E.~Perez, A.~Piktus \emph{et al.},
``Retrieval-augmented generation for knowledge-intensive NLP tasks,''
in \emph{Adv.\ Neural Inf.\ Process.\ Syst.}, vol.~33,
pp.~9459--9474, 2020.

\bibitem{gao2023rag}
Y.~Gao, Y.~Xiong, X.~Gao \emph{et al.},
``Retrieval-augmented generation for large language models: A survey,''
arXiv:2312.10997, 2023.

\bibitem{asai2023selfrag}
A.~Asai, Z.~Wu, Y.~Wang, A.~Sil, and H.~Hajishirzi,
``Self-RAG: Learning to retrieve, generate, and critique through
self-reflection,''
arXiv:2310.11511, 2023.

\bibitem{yang2020finbert}
Y.~Yang, M.~C.~S.~Uy, and A.~Huang,
``FinBERT: A pretrained language model for financial communications,''
arXiv:2006.08097, 2020.

\bibitem{devlin2019bert}
J.~Devlin, M.~W.~Chang, K.~Lee, and K.~Toutanova,
``BERT: Pre-training of deep bidirectional transformers for language
understanding,''
in \emph{Proc.\ NAACL-HLT}, pp.~4171--4186, 2019.

\bibitem{yao2023react}
S.~Yao, J.~Zhao, D.~Yu \emph{et al.},
``ReAct: Synergizing reasoning and acting in language models,''
in \emph{Proc.\ ICLR}, 2023.

\bibitem{langgraph2023}
LangChain,
``LangGraph: Build stateful, multi-actor applications with LLMs,''
2023. [Online]. Available: \url{https://langchain-ai.github.io/langgraph}

\bibitem{mcp2024}
Anthropic,
``Model Context Protocol specification,''
2024. [Online]. Available: \url{https://modelcontextprotocol.io}

\bibitem{newman2015}
S.~Newman,
\emph{Building Microservices: Designing Fine-Grained Systems}.
O'Reilly Media, 2015.

\bibitem{morganstanley2023}
Morgan Stanley,
``AI @ Morgan Stanley: Technology and innovation report,''
Internal Publication, 2023.

\bibitem{wu2023bloomberggpt}
S.~Wu, O.~Irsoy, S.~Lu \emph{et al.},
``BloombergGPT: A large language model for finance,''
arXiv:2303.17564, 2023.

\bibitem{yang2023fingpt}
H.~Yang, X.~Y.~Liu, and C.~D.~Wang,
``FinGPT: Open-source financial large language models,''
arXiv:2306.06031, 2023.

\bibitem{hutto2014vader}
C.~J.~Hutto and E.~Gilbert,
``VADER: A parsimonious rule-based model for sentiment analysis of
social media text,''
in \emph{Proc.\ ICWSM}, 2014.

\bibitem{wang2020minilm}
W.~Wang, F.~Wei, L.~Dong \emph{et al.},
``MiniLM: Deep self-attention distillation for task-agnostic compression
of pre-trained transformers,''
in \emph{Adv.\ Neural Inf.\ Process.\ Syst.}, vol.~33, 2020.

\bibitem{malo2014phrasebank}
P.~Malo, A.~Sinha, P.~Korhonen, J.~Wallenius, and P.~Takala,
``Good debt or bad debt: Detecting semantic orientations in economic
texts,''
\emph{J.~Assoc.\ Inf.\ Sci.\ Technol.}, vol.~65, no.~4,
pp.~782--796, 2014.

\bibitem{shuster2021}
K.~Shuster, S.~Poff, M.~Chen, D.~Kiela, and J.~Weston,
``Retrieval augmentation reduces hallucination in conversation,''
in \emph{Findings of EMNLP}, 2021.

\bibitem{karpas2022mrkl}
E.~Karpas, O.~Abend, Y.~Belinkov \emph{et al.},
``MRKL systems: A modular, neuro-symbolic architecture combining large
language models, external knowledge sources and discrete reasoning,''
arXiv:2205.00445, 2022.

\end{thebibliography}

\end{document}
